\documentclass[UTF8,11pt]{ctexart}

\usepackage[a4paper,margin=2.4cm]{geometry}
\usepackage{amsmath,amssymb,amsthm}
\usepackage{xcolor}
\usepackage{enumitem}
\usepackage{fancyhdr}
\usepackage{hyperref}

\definecolor{courseblue}{RGB}{43,49,139}
\hypersetup{colorlinks=true,linkcolor=courseblue,urlcolor=courseblue}
\setlength{\parindent}{2em}
\setlength{\parskip}{0.35em}
\setlist{itemsep=0.25em,topsep=0.35em}

\pagestyle{fancy}
\setlength{\headheight}{14pt}
\fancyhf{}
\lhead{复变函数补充阅读}
\rhead{02}
\cfoot{\thepage}
\renewcommand{\headrulewidth}{0.4pt}

\newcommand{\R}{\mathbb{R}}
\newcommand{\abs}[1]{\left\lvert#1\right\rvert}
\newcommand{\norm}[1]{\left\lVert#1\right\rVert}

\title{\textcolor{courseblue}{补充阅读 02：二元函数的可微性与增量形式}}
\author{}
\date{}

\begin{document}
\maketitle
\vspace{-2.2em}

\section{为什么复变函数会用到二元函数的可微性}

复变函数
\[
f(z)=u(x,y)+iv(x,y),\qquad z=x+iy,
\]
可以看作从 $\R^2$ 到 $\R^2$ 的映射
\[
(x,y)\longmapsto (u(x,y),v(x,y)).
\]
因此，在研究复函数 $f$ 的可导性之前，需要明确实二元函数
$u(x,y)$ 和 $v(x,y)$ 的可微性。这里的“可微”不是“两个偏导数存在”的简称，而是要求函数在该点附近能够被一个线性函数统一地近似。

\section{二元函数可微的定义}

设 $g(x,y)$ 定义在 $(x_0,y_0)$ 的某个邻域内。令
\[
\Delta g
=g(x_0+\Delta x,y_0+\Delta y)-g(x_0,y_0),
\qquad
r=\sqrt{(\Delta x)^2+(\Delta y)^2}.
\]
如果存在与 $\Delta x,\Delta y$ 无关的常数 $A,B$，使得
\[
\Delta g=A\Delta x+B\Delta y+R(\Delta x,\Delta y),
\]
并且
\[
\lim_{(\Delta x,\Delta y)\to(0,0)}
\frac{R(\Delta x,\Delta y)}{r}=0,
\]
那么称 $g$ 在 $(x_0,y_0)$ \textbf{可微}。

余项也常写成 $o(r)$，于是定义可以简写为
\[
\boxed{\Delta g=A\Delta x+B\Delta y+o(r).}
\]
其中 $A\Delta x+B\Delta y$ 是增量 $\Delta g$ 的\textbf{线性主部}，也就是函数在该点的一阶近似。

上述极限的含义是：对任意 $\varepsilon>0$，存在 $\delta>0$，使得
\[
0<r<\delta
\quad\Longrightarrow\quad
\abs{R(\Delta x,\Delta y)}\leq\varepsilon r.
\]
这个估计对所有趋近方向同时成立，因而比只沿坐标轴考察函数更强。

\section{线性主部为什么由偏导数给出}

如果 $g$ 在 $(x_0,y_0)$ 可微，那么取 $\Delta y=0$ 可得
\[
\frac{g(x_0+\Delta x,y_0)-g(x_0,y_0)}{\Delta x}
=A+\frac{R(\Delta x,0)}{\Delta x}.
\]
由于此时 $r=\abs{\Delta x}$，余项之比趋于 $0$，所以
\[
A=g_x(x_0,y_0).
\]
同理，取 $\Delta x=0$ 得到
\[
B=g_y(x_0,y_0).
\]
因此，二元函数可微时必有
\[
\boxed{
\Delta g
=g_x(x_0,y_0)\Delta x
+g_y(x_0,y_0)\Delta y
+o(r).
}
\]
其线性主部称为全微分：
\[
\mathrm{d}g=g_x\,\mathrm{d}x+g_y\,\mathrm{d}y.
\]

需要注意，以上结论说明“可微必然推出偏导数存在”，反过来仅有两个偏导数存在，一般不能推出可微。

\section{常用判断方法}

\subsection*{方法一：直接验证定义}

先找出候选线性主部
\[
L(\Delta x,\Delta y)=A\Delta x+B\Delta y,
\]
再证明
\[
\frac{\abs{\Delta g-L(\Delta x,\Delta y)}}{
\sqrt{(\Delta x)^2+(\Delta y)^2}}
\longrightarrow0.
\]
这一方法最接近定义，适合处理偏导数不连续或表达式中含有绝对值、根式的函数。

\subsection*{方法二：使用连续偏导数判据}

如果 $g_x$ 和 $g_y$ 在 $(x_0,y_0)$ 的某个邻域内存在，并且在
$(x_0,y_0)$ 连续，那么 $g$ 在该点可微。

这是常用的\textbf{充分条件}。对多项式、指数函数、三角函数及其通常的复合运算，往往可以直接使用这个判据。偏导数连续不是可微的必要条件，因此不能把它当作可微的定义。

\subsection*{判断关系}

在一点处，下列逻辑关系需要区分：
\[
\text{一阶偏导数在邻域存在且在该点连续}
\Longrightarrow
\text{可微}
\Longrightarrow
\begin{cases}
\text{连续},\\
\text{一阶偏导数存在}.
\end{cases}
\]
两个反方向一般都不成立。

\section{例子}

\subsection*{例 1：用连续偏导数判据}

设
\[
g(x,y)=x^2+xy+3y^2.
\]
其偏导数
\[
g_x=2x+y,\qquad g_y=x+6y
\]
在整个平面连续，所以 $g$ 在 $\R^2$ 上处处可微。在
$(x_0,y_0)$ 处，
\[
\Delta g
=(2x_0+y_0)\Delta x+(x_0+6y_0)\Delta y
+o(r).
\]

\subsection*{例 2：偏导数存在但不可微}

定义
\[
g(x,y)=
\begin{cases}
\sqrt{\abs{xy}},&(x,y)\neq(0,0),\\
0,&(x,y)=(0,0).
\end{cases}
\]
沿两条坐标轴都有 $g=0$，所以
\[
g_x(0,0)=g_y(0,0)=0.
\]
如果 $g$ 在原点可微，线性主部只能是 $0$，因而必须有
\[
\frac{g(\Delta x,\Delta y)}
{\sqrt{(\Delta x)^2+(\Delta y)^2}}\longrightarrow0.
\]
但沿 $\Delta y=\Delta x\neq0$，
\[
\frac{g(\Delta x,\Delta x)}
{\sqrt{2(\Delta x)^2}}
=\frac{\abs{\Delta x}}{\sqrt2\abs{\Delta x}}
=\frac1{\sqrt2}.
\]
极限不为 $0$，所以 $g$ 在原点不可微。这个例子表明，仅仅计算出偏导数还不够。

\section{与 Chapter 2 第 22 页证明的连接}

在课件中，假设复函数 $f$ 在 $z=x+iy$ 处可导，并写成
\[
f'(z)=a+ib,\qquad
\rho(\Delta z)=\rho_1+i\rho_2,
\qquad \rho(\Delta z)\to0.
\]
比较实部和虚部得到
\[
\Delta u
=a\Delta x-b\Delta y
+\rho_1\Delta x-\rho_2\Delta y,
\]
\[
\Delta v
=b\Delta x+a\Delta y
+\rho_2\Delta x+\rho_1\Delta y.
\]
下面用两种方法估计 $u$ 的余项。

\noindent\textbf{方法一：Cauchy--Schwarz 不等式。}
对于
\[
\mathbf{a}=(A,B),\qquad
\mathbf{b}=(X,Y),\qquad
\mathbf{a}\cdot\mathbf{b}=AX+BY,
\]
其中内积 $\mathbf{a}\cdot\mathbf{b}$ 就是对应分量相乘后相加。
Cauchy--Schwarz 不等式及其展开形式为
\[
\begin{aligned}
\boxed{\abs{\mathbf{a}\cdot\mathbf{b}}
       \leq\norm{\mathbf{a}}\,\norm{\mathbf{b}}}
\quad\Longleftrightarrow\quad
\abs{AX+BY}
&\leq\sqrt{A^2+B^2}\sqrt{X^2+Y^2}.
\end{aligned}
\]
在这里取
\[
\mathbf{a}=(\rho_1,-\rho_2),
\qquad
\mathbf{b}=(\Delta x,\Delta y),
\qquad
\mathbf{a}\cdot\mathbf{b}
=\rho_1\Delta x-\rho_2\Delta y,
\]
便得到
\[
\abs{\rho_1\Delta x-\rho_2\Delta y}
\leq
\sqrt{\rho_1^2+\rho_2^2}
\sqrt{(\Delta x)^2+(\Delta y)^2}
=\abs{\rho(\Delta z)}r=o(r).
\]

\noindent\textbf{方法二：三角不等式。}
由于 $\abs{\Delta x}\leq r$，$\abs{\Delta y}\leq r$，
\[
\begin{aligned}
\abs{\rho_1\Delta x-\rho_2\Delta y}
&\leq\abs{\rho_1}\abs{\Delta x}
     +\abs{\rho_2}\abs{\Delta y}\\
&\leq(\abs{\rho_1}+\abs{\rho_2})r=o(r).
\end{aligned}
\]
最后一步使用了 $\rho_1,\rho_2\to0$。

对 $v$ 的余项，Cauchy--Schwarz 不等式给出
\[
\abs{\rho_2\Delta x+\rho_1\Delta y}
\leq\abs{\rho(\Delta z)}r=o(r).
\]
三角不等式则给出
\[
\abs{\rho_2\Delta x+\rho_1\Delta y}
\leq(\abs{\rho_1}+\abs{\rho_2})r=o(r).
\]
因此这两个增量都具有“线性主部加 $o(r)$”的形式，从而 $u,v$
在 $(x,y)$ 可微，并且
\[
u_x=a,\qquad u_y=-b,\qquad
v_x=b,\qquad v_y=a.
\]
于是立即得到 Cauchy--Riemann 方程
\[
u_x=v_y,\qquad u_y=-v_x.
\]

\section{自测}

\begin{enumerate}
  \item 写出 $g(x,y)=e^x\cos y$ 在 $(0,0)$ 处的线性主部。
  \item 若
        \[
        \Delta g=2\Delta x-3\Delta y
        +(\Delta x)^2+\Delta x\Delta y,
        \]
        说明 $g$ 为什么可微，并求 $g_x,g_y$。
\end{enumerate}

\subsection*{参考答案}

\begin{enumerate}
  \item $\Delta g=\Delta x+o(r)$，线性主部为 $\Delta x$。
  \item 余项的绝对值不超过 $r^2+r^2/2$，所以余项为
        $o(r)$；$g_x=2$，$g_y=-3$。
\end{enumerate}

\end{document}
